% Generated by publication/build_sections.py; do not edit by hand.
\paragraph{Informal verdict.}
\textbf{A --- Complete solution}, with \textbf{high confidence}.

\paragraph{Lean coverage.}
Lean proves the unconditional first relative correction for Bell numbers, including the missing saddle-window and tail estimates and exact normalization.

\subsection{Audited informal answer}
\subsubsection{Faithful restatement}

Let \(T(n)\) be the \(n\)-th Bell number, written in Dobinski form as

\[
T(n)=e^{-1}\sum_{k=0}^{\infty}\frac{k^n}{k!}.
\]

Let \(w:[0,\infty\)\textbackslash{}to[1,\textbackslash{}infty)) be the inverse of \(y\mapsto y\log y\); thus

\[
w(x)\log w(x)=x.
\]

The earlier response R506 established

\[
T(n)\sim f(n),
\qquad
f(x)=\frac{w(x)^x e^{,w(x)-x-1}}{\sqrt{\log x}}.
\]

Q565 asks for the next term in this asymptotic expansion. There are no subparts or footnotes. There is no apparent mathematical typo; the only implicit domain point is that the displayed real-valued formula for \(f(x)\) requires \(x>1\), which is harmless in an asymptotic question as \(n\to\infty\). 

Throughout, \(\log\) denotes the natural logarithm.

\medskip\hrule\medskip

\subsubsection{Main result}

Set

\[
L=\log n,\qquad \ell=\log\log n.
\]

Then, as \(n\to\infty\) through the positive integers,

\[
\boxed{
T(n)
=
%
f(n)\left[
1+\frac{\ell-1}{2L}
+\frac{3\ell^2-10\ell+3}{8L^2}
+O\!\left(\frac{\ell^3}{L^3}\right)
\right].
}
\]

In particular, the requested next term is

\[
\boxed{
T(n)
=
%
f(n)
+
\frac{\log\log n-1}{2\log n},f(n)
+
O\!\left(
f(n)\frac{(\log\log n)^2}{(\log n)^2}
\right).
}
\]

Consequently,

\[
\boxed{
T(n)-f(n)
\sim
\frac{\log\log n}{2\log n},f(n).
}
\]

The expression with \(\log\log n-1\) is the more accurate first correction; the last display gives its simpler asymptotic equivalent.

A stronger saddle-point form is also valid. If

\[
r=\log w(n),
\]

so that \(re^r=n\) and \(w(n)=e^r=n/r\), then

\[
\boxed{
T(n)
=
%
\frac{w(n)^n e^{,w(n)-n-1}}{\sqrt{1+r}}
\left[
1-
\frac{r^2(2r^2+7r+10)}
{24n(r+1)^3}
+
O\!\left(\left(\frac rn\right)^{3/2}\right)
\right].
}
\]

All \(O\)-constants here are absolute, and the estimates hold uniformly as \(n\to\infty\) through the integers.

\medskip\hrule\medskip

\subsection{Proof}

\subsubsection{1. Exponential generating function}

Write \(B_n=T(n)\). From Dobinski's formula,

\[
\begin{aligned}
\sum_{n=0}^{\infty}B_n\frac{z^n}{n!}
&=
e^{-1}\sum_{k=0}^{\infty}\frac1{k!}
\sum_{n=0}^{\infty}\frac{(kz)^n}{n!}  \\
&=
e^{-1}\sum_{k=0}^{\infty}\frac{e^{kz}}{k!}
=
%
= e^{-1}e^{e^z}
%
\exp(e^z-1).
\end{aligned}
\]

The interchange is justified by absolute and locally uniform convergence. Hence

\[
\frac{B_n}{n!}=[z^n]\exp(e^z-1).
\]

Let \(r>0\) be the unique solution of

\[
re^r=n.
\]

This is the saddle-point equation. Put

\[
b=n(r+1).
\]

Cauchy's coefficient formula on the circle \(|z|=r\) gives

\[
\frac{B_n}{n!}
=
%
\frac{e^{e^r-1}}{2\pi r^n}J_n,
\]

where

\[
J_n
=
%
\int_{-\pi}^{\pi}e^{\Phi(\theta)},d\theta,
\qquad
\Phi(\theta)
=
%
e^{r e^{i\theta}}-e^r-in\theta.
\]

The choice \(re^r=n\) ensures \(\Phi'(0)=0\).

\medskip\hrule\medskip

\subsubsection{2. Negligibility of the noncentral arc}

For \(-\pi\leq\theta\leq\pi\),

\[
\begin{aligned}
\Re\Phi(\theta)
&=
e^{r\cos\theta}\cos(r\sin\theta)-e^r\\
&\leq e^{r\cos\theta}-e^r\\
&=
-e^r\left(1-e^{-r(1-\cos\theta)}\right).
\end{aligned}
\]

There is an absolute \(c>0\) such that

\[
1-e^{-x}\geq c\min(x,1)\qquad(x\geq0),
\]

and

\[
1-\cos\theta\geq \frac{2\theta^2}{\pi^2}
\qquad (|\theta|\leq\pi).
\]

Since \(e^r=n/r\), it follows that

\[
\Re\Phi(\theta)
\leq
-c\min(n\theta^2,n/r).
\tag{1}
\]

Define

\[
q=\frac rn=e^{-r}=\frac1{w(n)}
\]

and choose

\[
T=q^{-1/20},
\qquad
\theta_0=\frac{T}{\sqrt b}.
\]

Then

\[
n\theta_0^2=\frac{T^2}{r+1}
=\frac{q^{-1/10}}{r+1}\longrightarrow\infty.
\]

Consequently, by \(1\),

\[
\int_{\theta_0\leq|\theta|\leq\pi}
\left\lvert{}e^{\Phi(\theta)}\right\rvert{},d\theta
\leq
2\pi
\exp\left(-c\frac{q^{-1/10}}{r+1}\right).
\tag{2}
\]

This is smaller than \(b^{-1/2}q^M\) for every fixed \(M>0\). Thus only the arc \(|\theta|\leq\theta_0\) contributes to any algebraic order in \(q\).

\medskip\hrule\medskip

\subsubsection{3. Expansion near the saddle}

Direct differentiation gives

\[
\Phi''(0)=-n(r+1)=-b,
\]

\[
\Phi'''(0)
=
%
-i,n(r^2+3r+1),
\]

and

\[
\Phi^{(4)}(0)
=
%
n(r^3+6r^2+7r+1).
\]

Moreover, on the central arc,

\[
|\Phi^{(5)}(\theta)|\leq Cnr^4
\]

for an absolute constant \(C\). Taylor's formula therefore yields

\[
\Phi(\theta)
=
%
-\frac b2\theta^2
-\frac{ic}{6}\theta^3
+\frac d{24}\theta^4
+O(nr^4|\theta|^5),
\tag{3}
\]

where

\[
c=n(r^2+3r+1),
\qquad
d=n(r^3+6r^2+7r+1).
\]

Set \(t=\sqrt b,\theta\). Then \(3\) becomes

\[
\Phi\left(\frac{t}{\sqrt b}\right)
=
%
-\frac{t^2}{2}
-i\lambda_3t^3
+\lambda_4t^4
+O(q^{3/2}|t|^5),
\tag{4}
\]

where

\[
\lambda_3=\frac{c}{6b^{3/2}}=O(q^{1/2}),
\qquad
\lambda_4=\frac{d}{24b^2}=O(q).
\]

For \(|t|\leq T=q^{-1/20}\), all the perturbation terms in \(4\) tend uniformly to zero. Expanding the exponential gives, for some fixed polynomial \(P\),

\[
\begin{aligned}
\exp\left(
-i\lambda_3t^3+\lambda_4t^4+O(q^{3/2}|t|^5)
\right)
={}&
1-i\lambda_3t^3+\lambda_4t^4
-\frac12\lambda_3^2t^6\\
&\quad+O\bigl(q^{3/2}P(|t|)\bigr).
\end{aligned}
\tag{5}
\]

The term in \(t^3\) integrates to zero over the symmetric interval. The standard Gaussian moments, obtained by integration by parts, are

\[
\frac1{\sqrt{2\pi}}
\int_{\mathbb R}t^4e^{-t^2/2},dt=3,
\qquad
\frac1{\sqrt{2\pi}}
\int_{\mathbb R}t^6e^{-t^2/2},dt=15.
\]

Using \(2\)--\(5\), extending the Gaussian integral from ([-T,T]) to \(\mathbb R\), and noting that the resulting tails are exponentially small, we obtain

\[
J_n
=
%
\sqrt{\frac{2\pi}{b}}
\left[
1+3\lambda_4-\frac{15}{2}\lambda_3^2
+O(q^{3/2})
\right].
\tag{6}
\]

Now

\[
3\lambda_4
=
%
\frac{r^3+6r^2+7r+1}
{8n(r+1)^2},
\]

whereas

\[
\frac{15}{2}\lambda_3^2
=
%
\frac{5(r^2+3r+1)^2}
{24n(r+1)^3}.
\]

Putting these over a common denominator gives

\[
3\lambda_4-\frac{15}{2}\lambda_3^2
=
%
-\frac{2r^4+9r^3+16r^2+6r+2}
{24n(r+1)^3}.
\tag{7}
\]

Consequently,

\[
\frac{B_n}{n!}
=
%
\frac{e^{e^r-1}}{r^n\sqrt{2\pi n(r+1)}}
\left[
1-
\frac{2r^4+9r^3+16r^2+6r+2}
{24n(r+1)^3}
+O(q^{3/2})
\right].
\tag{8}
\]

\medskip\hrule\medskip

\subsubsection{4. Incorporating Stirling's expansion}

We use Stirling's formula with its first correction:

\[
n!
=
%
\sqrt{2\pi n}\left(\frac ne\right)^n
\left[
1+\frac1{12n}+O(n^{-2})
\right].
\tag{9}
\]

For completeness, \(9\) follows from the Euler--Maclaurin expansion

\[
\log n!
=
%
\left(n+\frac12\right)\log n-n
+\frac12\log(2\pi)
+\frac1{12n}
+O(n^{-3});
\]

the constant is identified by Wallis's product, and exponentiation yields \(9\).

Multiplying \(8\) by \(9\), and using \(n/r=e^r=w(n)\), gives

\[
B_n
=
%
\frac{w(n)^n e^{w(n)-n-1}}{\sqrt{r+1}}
\left[
1+\frac1{12n}
-\frac{2r^4+9r^3+16r^2+6r+2}
{24n(r+1)^3}
+O(q^{3/2})
\right].
\]

The two displayed correction terms combine as follows:

\[
\begin{aligned}
\frac1{12n}
&-
\frac{2r^4+9r^3+16r^2+6r+2}
{24n(r+1)^3}\\
&=
-\frac{r^2(2r^2+7r+10)}
{24n(r+1)^3}.
\end{aligned}
\]

Thus

\[
B_n
=
%
\frac{w(n)^n e^{w(n)-n-1}}{\sqrt{r+1}}
\left[
1-
\frac{r^2(2r^2+7r+10)}
{24n(r+1)^3}
+O\!\left(q^{3/2}\right)
\right],
\]

which proves the stronger saddle-point formula.

\medskip\hrule\medskip


Recall that

\[
f(n)=
\frac{w(n)^n e^{w(n)-n-1}}{\sqrt L},
\qquad L=\log n.
\]

Therefore

\[
\frac{B_n}{f(n)}
=
%
\sqrt{\frac L{r+1}}
\left[
1+O\left(\frac rn\right)
\right].
\tag{10}
\]

The \(O(r/n)\) term is smaller than every fixed negative power of \(\log n\), so the logarithmic correction comes entirely from the factor \(\sqrt{L/(r+1)}\).

Since \(re^r=n\), taking logarithms gives

\[
r+\log r=L.
\tag{11}
\]

Let \(\ell=\log L\). We next expand the solution of \(11\).

Because \(r<L\) and \(\log r\leq\log L=\ell\),

\[
L-\ell\leq r<L.
\]

Write

\[
r=L-\ell+\delta.
\]

Then \(0\leq\delta\leq\ell\), and \(11\) gives

\[
\delta+
\log\left(1-\frac{\ell-\delta}{L}\right)=0.
\]

Since \(0\leq(\ell-\delta)/L\leq\ell/L\to0\),

\[
\delta
=
%
\frac{\ell-\delta}{L}
+
O\left(\frac{\ell^2}{L^2}\right).
\]

Hence

\[
\delta
=
%
\frac{\ell}{L}
+
O\left(\frac{\ell^2}{L^2}\right),
\]

and therefore

\[
r
=
%
L-\ell+\frac{\ell}{L}
+
O\left(\frac{\ell^2}{L^2}\right).
\tag{12}
\]

It follows that

\[
\frac{r+1}{L}
=
%
1+\frac{1-\ell}{L}
+\frac{\ell}{L^2}
+
O\left(\frac{\ell^2}{L^3}\right).
\tag{13}
\]

Apply

\[
(1+u)^{-1/2}
=
%
1-\frac u2+\frac{3u^2}{8}+O(u^3)
\]

to the right-hand side of \(13\). Since \(u=O(\ell/L)\), we find

\[
\sqrt{\frac L{r+1}}
=
%
1+\frac{\ell-1}{2L}
+\frac{3\ell^2-10\ell+3}{8L^2}
+
O\left(\frac{\ell^3}{L^3}\right).
\tag{14}
\]

Finally, \(r/n=O(L/n)\), so the saddle correction in \(10\) is

\[
O(L/n)
=
%
o\left(\frac{\ell^3}{L^3}\right).
\]

Combining \(10\) and \(14\) proves

\[
T(n)
=
%
f(n)\left[
1+\frac{\ell-1}{2L}
+\frac{3\ell^2-10\ell+3}{8L^2}
+
O\left(\frac{\ell^3}{L^3}\right)
\right].
\]

This completes the proof.

\medskip\hrule\medskip

\subsubsection{Boundary cases and numerical checks}

The asymptotic statement concerns integers \(n\to\infty\). The function \(f(n)\) is not finite at \(n=1\), because its denominator is \(\sqrt{\log 1}=0\); thus no assertion is made there. The Bell numbers themselves satisfy \(T(0)=T(1)=1\), with the usual convention \(0^0=1\) in Dobinski's formula.

For sufficiently large \(n\), \(\ell-1>0\), and the error is smaller than the first correction. Hence the expansion also proves

\[
T(n)>f(n)
\]

for all sufficiently large \(n\). No such finite-\(n\) inequality is implicit in the asymptotic statement.

The following values are an independent numerical check, not part of the proof. Here

\[
R_1=1+\frac{\ell-1}{2L},
\qquad
R_2=R_1+\frac{3\ell^2-10\ell+3}{8L^2}.
\]

\[
\begin{array}{c|c|c|c}
n&T(n)/f(n)&R_1&R_2\\ \hline
20&0.956182538&1.016221193&0.955487159\\
100&1.021429523&1.057237801&1.026146557\\
1000&1.050833863&1.067507077&1.054091807
\end{array}
\]

The convergence is logarithmically slow; the \(L^{-2}\) term is therefore still significant at moderate \(n\).

A reproducible Python check is:

\begin{verbatim}
from math import comb
import mpmath as mp

mp.mp.dps = 80


def bell_numbers(n_max: int) -> list[int]:
    """Return B_0, ..., B_n_max using B_{m+1} = sum_k binom(m,k) B_k."""
    values = [1]
    for m in range(n_max):
        values.append(sum(comb(m, k) * values[k] for k in range(m + 1)))
    return values


def check(n_values: list[int]) -> None:
    bells = bell_numbers(max(n_values))

    print("n        T(n)/f(n)              R1                      R2")
    for n in n_values:
        L = mp.log(n)
        ell = mp.log(L)
        r = mp.re(mp.lambertw(n))
        w = n / r

        log_f = n * mp.log(w) + w - n - 1 - mp.log(L) / 2
        exact_ratio = mp.mpf(bells[n]) * mp.exp(-log_f)

        R1 = 1 + (ell - 1) / (2 * L)
        R2 = R1 + (3 * ell**2 - 10 * ell + 3) / (8 * L**2)

        print(
            f"{n:<8d} "
            f"{mp.nstr(exact_ratio, 14):<24} "
            f"{mp.nstr(R1, 14):<24} "
            f"{mp.nstr(R2, 14)}"
        )


check([20, 100, 1000])
\end{verbatim}

Thus the precise first correction requested by Q565 is

\[
\boxed{
\frac{\log\log n-1}{2\log n},f(n),
}
\]

with relative remainder \(O((\log\log n)^2/(\log n)^2)\).
